% p1-12-flos-aureus-link

\section{P1 §12. Relationship to the Trinity $S^{3}$AI / Flos Aureus PhD}

12. Relationship to the Trinity $S^{3}$AI / Flos Aureus PhD
  Monograph and Stergios Collaboration Note

         Figure (fig-p1-ch12-monograph). Provenance triptych --- Monograph spine / Shared anchor
                                   identity / Collaboration handshake.

  12.1 Scope of the Relationship
  This paper extracts a narrow, journal-safe statistical subproblem from the broader
  Trinity $S^{3}$AI / Flos Aureus research program, whose PhD monograph is in preparation
  under the title Trinity $S^{3}$AI --- Flos Aureus v6.2 / GOLDEN SUNFLOWERS (D. Vasilev,
  ORCID: 0009-0008-4294-6159, defense scheduled 2026-06-15). The relationship
  between this paper and the monograph is one of provenance, algebraic foundation, and
  falsification design --- not one of external validation. The monograph does not
  independently validate the statistical claims made in this paper; those claims stand or
  fall entirely on the null-model analysis and MDL scoring defined in Sections 5--7.

  Specifically:
  - The monograph is not cited as evidence that G$\varphi$ achieves non-random compressibility.
  - The monograph's exploratory catalogue (approximately 42 fits, Chapter 34) is treated
  as historical motivation and a supplementary candidate pool; expressions from it are
  re-evaluated under the null-model analysis.
  - The identity $\varphi$2 + $\varphi$-2 = 3 (PhD: Ch.3 Trinity Identity) is used only as an algebraic
  completeness statement about the Lucas ring Z[$\varphi$].

  12.2 Monograph Contributions Table

  12.3 Proof Status and Honest Admission
  The PhD monograph maintains a Coq citation map and an R5-honest admitted-budget
  ledger (PhD: App.F Coq Map). A local source audit of the available t27 checkout gives
  the following status. The topic-specific physical-constant layer t27/proofs/trinity/
  contains 13 Coq files, 281 theorem-like declarations, 122 source-level Qed. markers, 32
  source-level Admitted. markers, and 0 Abort. markers. The broader t27/proofs/ tree

  contains 18 Coq proof files, 296 theorem-like declarations, 127 source-level Qed.
  markers, 32 source-level Admitted. markers, and 0 Abort. markers. The older t27/coq/
  kernel layer contains 10 Coq files, 73 theorem-like declarations, 35 source-level Qed.
  markers, 0 Admitted. markers, and 11 source-level Abort. markers. Catalog42.v itself
  contains 13 source-level Qed. markers and 0 Admitted. markers, and defines a
  totalverifiedtheorems aggregate over 28 representative theorem obligations across
  core identities, AlphaPhi, gauge, CKM, neutrino, mass, and monomial-interface claims.
  Therefore the correct statement is: the paper has a 42-row formula catalogue, a
  69-formula evaluator interface noted in FormulaEval.v, and a larger Coq proof layer
  with 122 Qed. markers in t27/proofs/trinity/; it is not yet a closed one-theorem-per-row
  proof set for all 42 or all 69 formula entries. Because coqc is not available in the current
  build sandbox, this is a source audit rather than a fresh local Coq compilation verdict.

  12.4 Stergios Collaboration Note
  The Pellis hierarchical expansion of Section 9 is named after Stergios Pellis, who is a
  co-author of this paper and the principal contributor of the multi-scale decomposition
  framework. The collaboration between Dmitrii Vasilev and Stergios Pellis originates
  within the Trinity $S^{3}$AI research program and represents the integration of two
  complementary methodological threads:

  Vasilev's thread is the algebraic and information-theoretic foundation: the formal
  grammar G$\varphi$, the MDL scoring framework, the multiple-testing correction architecture,
  and the PhD monograph's algebraic backbone (the Lucas ring, the Trinity identity $\varphi$2 +
  $\varphi$-2 = 3, and the Flos Aureus numerical analysis). This thread focuses on rigor,
  falsifiability, and statistical correctness.

  Pellis's thread is the hierarchical expansion methodology: the insight that a
  multiplicative grammar can be systematically refined by additive corrections in powers
  of $\varphi$-1, defining a filtration F1 ⊃ F2 ⊃ … over which constants admit multi-scale
  decompositions. The Pellis expansion has independent mathematical interest as a
  structured approximation theory over the ring Z[$\varphi$], connecting to generalized
  continued fraction expansions and to the theory of beta-expansions in non-integer
  bases.

  The collaboration is ongoing. The current paper represents the first joint formalization
  of the combined framework. Future work under this collaboration includes:
  - Systematic comparison of the Pellis expansion against standard continued fractions
  and Engel expansions as approximation schemes for transcendental numbers.
  - Application of the Pellis filtration to near-threshold constants where the base grammar
  achieves only O(10-2) precision.
  - A dedicated paper (Paper 3 in the series) formalizing the hierarchical expansion and
  its MDL theory.

  The collaboration note is included here to establish clear provenance for the Pellis
  expansion methodology and to acknowledge its independent conceptual origin within
  the joint research program.

  12.5 Scott Olsen Historical-Geometric Note

  Scott A. Olsen's golden-section work is used here as Tier D context: it helps explain why
  the reader should expect a long historical lineage around $\varphi$, but it does not validate any
  numerical claim in this paper. Olsen's book The Golden Section: Nature's Greatest
  Secret presents the golden section as a unifying theme across mathematics, art,
  architecture, music, philosophy, and natural form (Internet Archive bibliographic
  record). That role is useful for exposition because the GOLDEN BRIDGE (Trinity-Pellis Bridge) program also uses $\varphi
  as a structural organizer, but the evidentiary standards are different.

  The specific Olsen contribution imported into this version is the golden balance motif:
  the ordered five-point reciprocal scale

  $\varphi$-2, $\varphi$-1, 1, $\varphi$, $\varphi$2.

  Within the article, this motif is treated as a geometric and pedagogical diagram of
  reciprocal scaling around the unit element. It is not used to infer any Standard Model
  constant, and it is not part of the null-model statistics. In PhD notation it belongs to the
  visual and historical layer, not to the proof ledger.

  This limited use protects the paper from a common reviewer objection. The history of
  the golden section is allowed to provide orientation, but every scientific claim remains
  inside the formal grammar, the explicit search space, the MDL score, and the
  falsification protocol.

  12.6 Author Contributions
  - Dmitrii Vasilev: Trinity monomial grammar, computational search framing,
  falsification protocol, reproducibility architecture, and integration with the Trinity $S^{3}$AI
  / Flos Aureus programme.
  - Stergios Pellis: Pellis hierarchical expansion framework, fine-structure comparison
  layer, mechanism-risk assessment, and mathematical co-development of the
  GOLDEN BRIDGE (Trinity-Pellis Bridge) series.
  - Scott Olsen: historical-geometric and golden-section context, golden-balance
  interpretive lineage, and visual-theoretical framing. These contributions are treated as
  Tier D context and are not used as statistical evidence, derivation evidence, or physical
  proof.

  12.7 What This Paper Does Not Inherit from the Monograph
  - This paper does not adopt the AlphaPhi derivation inconsistency in Chapter 38.
  - This paper does not use the ~42 catalogue expressions as pre-registered predictions.
  - This paper does not claim that Zenodo DOIs associated with the monograph constitute
  proof of statistical claims.
  - This paper does not import any stale proof-count claims; Coq obligation counts must
  be regenerated from the repository at submission time.

\subsection{References}

- Vasilev, D. \textit{Flos Aureus Monograph: Trinity Constants}. DOI \href{https://zenodo.org/doi/10.5281/zenodo.19227877}{10.5281/zenodo.19227877}.
- Pellis, S. "The Fine-Structure Constant and the Golden Ratio." \href{https://vixra.org/pdf/2110.0117v4.pdf}{viXra 2110.0117 v4}.
- Sherbon, M. A. "Fundamental Nature of the Fine-Structure Constant." \href{https://hal.science/hal-02341850/document}{HAL hal-02341850}.
